\begin{Lemma} \label{lemma:mainpair}
	Let $m \geq 1$ and $0 \leq a \leq m$. Suppose that $(\alpha_n,\beta_n)$ is a Bailey pair relative to $q$ and that there is a sequence $(f_n)_{n \geq -1}$ such that
	\begin{align} \label{eqn:fdeflater}
		\alpha_n=\dfrac{1-q^{2n+1}}{1-q}f_n,
	\end{align}
	where
	\begin{align}
		q^{-1}f_{-1}=-f_0.
		\label{eqn:f-1f0}
	\end{align}
	Then $(\alpha_n',\beta_n')$ is a Bailey pair relative to $q$, where
	\begin{equation*} %\label{eqn:mainpairalpha}
		\alpha_n' =
		\begin{cases}
			\dfrac{1}{1-q}q^{m\binom{n+1}{2}}\bigl(
			q^{a(n+1)}f_{n+1} -q^{-an+2n-1}f_{n-1}
			\bigr) & \text{if $0 \leq a \leq m-1$}, \vspace{1mm}\\
			\dfrac{1}{1-q}q^{m \binom{n+1}{2}}\bigl(1-q^{2n+1}\bigr)f_n &\text{if $a=m$},
		\end{cases}
	\end{equation*}
	and
	\begin{align}
		\beta_n' &{}= \beta_{n_{m+1}}' \nonumber \\
 &{}= \frac{1}{(-q)_{n_{m+1}}} \sum_{n_m, \dots, n_1 \geq 0} \frac{q^{\binom{n_m+1}{2} + \cdots + \binom{n_1+1}{2}}\bigl(q^2;q^2\bigr)_{n_1+\delta_{a,0}}}{(q)_{n_{m+1}-n_m} (q)_{n_m}} \beta_{n_1+\delta_{a,0}} \prod_{i=1}^{m-1} \begin{bmatrix} n_{i+1} + \delta_{a,i} \\ n_i \end{bmatrix}.\label{eqn:mainpairbeta}
	\end{align}
	
\end{Lemma}

\begin{proof}
	{\it Step 1\emph{:}} First assume $a > 0$. Iterate Lemma \ref{Baileylemma} $a$ times with $x\mapsto q$, $b\mapsto -q$, and $c\rightarrow \infty$ to obtain another Bailey pair relative to $q$. The resulting $\alpha_n$ and $\beta_n$ are
	\begin{align}
		\alpha_n=q^{a\binom{n+1}{2}}\dfrac{1-q^{2n+1}}{1-q}f_n
		\label{eqn:alphafn}
	\end{align}
	and
	\begin{align}
		\beta_n=\beta_{n_{a+1}}=
		\dfrac{1}{(-q)_{n_{a+1}}}\sum_{n_a,\dots,n_1\geq 0}
		\dfrac{q^{\binom{n_a+1}{2}+\cdots+\binom{n_1+1}{2}}(-q)_{n_1}}{(q)_{n_{a+1}-n_a}\cdots (q)_{n_2-n_1}}\beta_{n_1}. \label{eqn:betastep1}
	\end{align}
	Here and throughout the paper we use the convention that $1/(q)_n = 0$ if $n <0$.
	If $a=m$, we stop here. Multiplying the numerator and denominator of the above sum by $(q)_{n_1} \cdots (q)_{n_m}$ and rewriting in terms of $q$-binomial coefficients using \eqref{eqn:qbinom} gives \eqref{eqn:mainpairbeta} for $a=m$.
	
	{\it Step 2\emph{:}} Apply Lemma \ref{littlelemma} with $x \mapsto q$ to obtain a Bailey pair relative to $q^2$:
	\begin{align*}
		\alpha_n &{}= \dfrac{1-q^{n+1}}{(q)_2}\bigl(
		q^{a\binom{n+2}{2}}f_{n+1} + q^{n+a\binom{n+1}{2}}f_n\bigr)\\
		&{}=\dfrac{1-q^{n+1}}{(q)_2}q^{(a+2)\binom{n+1}{2}-n^2}
		\bigl(
		f_n+q^{a(n+1)-n}f_{n+1}\bigr)
	\end{align*}
	and
	\begin{align*}
		\beta_n=\beta_{n_{a+1}}=
		\dfrac{1-q^{n_{a+1}+1}}{(-q)_{n_{a+1}+1}}\sum_{n_a,\dots,n_1\geq 0}
		\dfrac{q^{\binom{n_a+1}{2}+\cdots+\binom{n_1+1}{2}}(-q)_{n_1}}{(q)_{n_{a+1}+1-n_a}\cdots (q)_{n_2-n_1}}\beta_{n_1}.
	\end{align*}
	
	
	{\it Step 3\emph{:}} Use Lemma \ref{Baileylemma} with $x\mapsto q^2$, $b\mapsto-q^2$, and $c \to \infty$ to obtain another Bailey pair relative to $q^2$. The resulting pair is
	\begin{align*}
		\alpha_n
		=\dfrac{1-q^{2n+2}}{(1+q)(q)_2}q^{(a+3)\binom{n+1}{2}-n^2}
		\bigl(
		f_n+q^{a(n+1)-n}f_{n+1}\bigr)
	\end{align*}
	and
	\begin{align*}
		\beta_n=\beta_{n_{a+2}}=
		\dfrac{1}{(1+q)(-q)_{n_{a+2}}}\sum_{n_{a+1},\dots,n_1\geq 0}
		\dfrac{q^{\binom{n_{a+1}+1}{2}+\cdots+\binom{n_1+1}{2}}\bigl(1-q^{n_{a+1}+1}\bigr)(-q)_{n_1}}{(q)_{n_{a+2}-n_{a+1}}(q)_{n_{a+1}+1-n_a}\cdots (q)_{n_2-n_1}}\beta_{n_1}.
	\end{align*}
	At this point, we multiply both
	$\alpha_n$ and $\beta_n$ by $1+q$.
	
	
	{\it Step 4\emph{:}} Use Lemma \ref{Baileylatticelemma} with $x \mapsto q^2$ and $b,c \mapsto q$ to obtain a Bailey pair relative to $q$. This retains the $\beta_n$, but does change the $\alpha_n$. When $n \geq 1$, we have
	\begin{gather*}
		\alpha_n = \dfrac{1}{1-q}
		\bigl(
		q^{(a+3)\binom{n+1}{2}-n^2}\bigl(f_n+q^{a(n+1)-n}f_{n+1}\bigr) \\ \hphantom{\alpha_n = \dfrac{1}{1-q}\bigl(}{}
		- q^{2n+(a+3)\binom{n}{2}-(n-1)^2}\bigl(f_{n-1}+q^{an-n+1}f_{n}\bigr)
		\bigr)\\ \hphantom{\alpha_n}{}
		=
		\dfrac{1}{1-q}
		\bigl(
		q^{(a+3)\binom{n+1}{2}-n^2+a(n+1)-n}f_{n+1}
		-
		q^{2n+(a+3)\binom{n}{2}-(n-1)^2}f_{n-1}
		\bigr)\\ \hphantom{\alpha_n}{}
		= \dfrac{1}{1-q}q^{(a+1)\binom{n+1}{2}}\bigl(
		q^{a(n+1)}f_{n+1} -q^{-an+2n-1}f_{n-1}
		\bigr).
	\end{gather*}
	When $n=0$, the above is also valid, using \eqref{eqn:latticeextracond} and \eqref{eqn:f-1f0}.
	
	
	{\it Step 5\emph{:}} Iterate Lemma \ref{Baileylemma} $m-a-1$ times with $x\mapsto q$, $b\mapsto -q$, and $c\rightarrow \infty$ to arrive at the final Bailey pair relative to $q$:
	\begin{align*}
		\alpha_n
		=
		\dfrac{1}{1-q}q^{m\binom{n+1}{2}}\bigl(
		q^{a(n+1)}f_{n+1} -q^{-an+2n-1}f_{n-1}
		\bigr)
	\end{align*}
	and
	\begin{align*}
		\beta_n=\beta_{n_{m+1}}=
		\dfrac{1}{(-q)_{n_{m+1}}}\sum_{n_m,\dots,n_1\geq 0}
		\dfrac{q^{\binom{n_{m}+1}{2}+\cdots+\binom{n_1+1}{2}}\bigl(1-q^{n_{a+1}+1}\bigr)(-q)_{n_1}}
		{(q)_{n_{m+1}-n_{m}}\cdots
			(q)_{n_{a+1}+1-n_a}\cdots (q)_{n_2-n_1}}\beta_{n_1}.
	\end{align*}
	Rewriting the sum in terms of $q$-binomial coefficients gives{\samepage
	\begin{equation*} %\label{eqn:genbetafin}
		\beta_n = \beta_{n_{m+1}} = \frac{1}{(-q)_{n_{m+1}}} \sum_{n_m, \dots, n_1 \geq 0} \frac{q^{\binom{n_m+1}{2} + \cdots + \binom{n_1+1}{2}}\bigl(q^2;q^2\bigr)_{n_1}}{(q)_{n_{m+1}-n_m} (q)_{n_m}} \beta_{n_1} \prod_{i=1}^{m-1} \begin{bmatrix} n_{i+1} + \delta_{a,i} \\ n_i \end{bmatrix},
	\end{equation*}
	which coincides with \eqref{eqn:mainpairbeta}.}
	
	{\it Step 6\emph{:}}
	If $a=0$, performing Steps 2--5 leads to the following $\beta_n$ (while the formula for $\alpha_n$ is the same as in Step 5 with $a\mapsto 0$):
	\begin{align*}
		\beta_n=\beta_{n_{m+1}}=
		\dfrac{1}{(-q)_{n_{m+1}}}\sum_{n_m,\dots,n_1\geq 0}
		\dfrac{q^{\binom{n_{m}+1}{2}+\cdots+\binom{n_1+1}{2}}\bigl(1-q^{n_{1}+1}\bigr)(-q)_{n_1+1}}
		{(q)_{n_{m+1}-n_{m}}\cdots (q)_{n_2-n_1}}
		\beta_{n_1+1}.
	\end{align*}
	Rewriting this in terms of $q$-binomial coefficients gives
	\begin{equation*} %\label{eqn:genbetafina0}
		\beta_n = \beta_{n_{m+1}} = \frac{1}{(-q)_{n_{m+1}}} \sum_{n_m, \dots, n_1 \geq 0} \frac{q^{\binom{n_m+1}{2} + \cdots + \binom{n_1+1}{2}}\bigl(q^2;q^2\bigr)_{n_1+1}}{(q)_{n_{m+1}-n_m} (q)_{n_m}} \beta_{n_1+1} \prod_{i=1}^{m-1} \begin{bmatrix} n_{i+1} + \delta_{a,i} \\ n_i \end{bmatrix},
	\end{equation*}
	which again matches \eqref{eqn:mainpairbeta}. This completes the proof.
\end{proof}
